This section introduces the background on static vulnerability checkers, LLM-based checker generation from security patches, and the analyzer-internal intermediate results that motivate our approach.

\subsection{Static Vulnerability Checkers}

Static analysis examines source code without executing it, aiming to find bugs, security vulnerabilities, and code-quality issues at compile time or during code review~\cite{engler2001bugs,bessey2010coverity}.
A vulnerability checker is a rule or query within a static-analysis framework that flags program locations matching a security-relevant condition.
Modern frameworks support different checker paradigms.
Path-sensitive engines such as the Clang Static Analyzer (CSA) perform symbolic exploration over program paths, tracking branch guards, symbolic values, memory regions, and lifetime transitions to detect bugs like use-after-free, null dereference, and resource leaks~\cite{clangsaDocs}.
Relational query engines such as CodeQL compile source code into a queryable database and express checkers as declarative queries over data-flow edges, call-graph relations, type hierarchies, and API-usage patterns~\cite{avgustinov2016ql,codeqlDocs}.
Both paradigms have been deployed at industrial scale for continuous vulnerability detection~\cite{sadowski2015tricorder,calcagno2015infer}.

Writing effective checkers remains a bottleneck.
A checker must encode not only the syntactic pattern of a vulnerability but also the semantic conditions (path guards, state invariants, data-flow relations, and API contracts) that distinguish real bugs from safe code.
This requires deep domain knowledge of both the target codebase and the analysis framework, making manual checker authoring expensive and error-prone.

\subsection{LLM-Based Patch-Driven Checker Generation}

Recent work has proposed using LLMs to automate checker generation from security patches.
A security patch provides a compact before/after specification: the vulnerable revision should trigger the generated checker, and the fixed revision should not.
This idea builds on a well-established line of work that treats patches as vulnerability specifications, including semantic patching~\cite{padioleau2008coccinelle}, clone-based vulnerability discovery~\cite{jang2012redebug,kim2017vuddy}, and patch-derived datasets~\cite{fan2020bigvul}.
KNighter~\cite{yang2025knighter} is a representative system that synthesizes CSA checkers from historical Linux kernel fixes and validates them against the vulnerable and fixed revisions.

However, the generated checkers are seldom semantically adequate.
A patch is inherently underspecified: a one-line edit may encode a path guard, an ownership transition, a lifetime barrier, or an API contract, yet the generated checker often captures only the syntactic shape of the edit without the semantic condition the edit enforces.
Existing refinement loops, including KNighter's own iterative strategy, operate at the output level~\cite{yang2025knighter,madaan2023selfrefine,shinn2023reflexion}.
They feed behavioral signals (e.g., ``the fixed side still warns'') back to the LLM, revealing the symptom but not diagnosing the cause.
This feedback cannot tell the LLM whether the missing predicate is a path guard, a cleanup barrier, an ownership transition, or an API contract, so the model must guess which semantic condition to add.

\subsection{Analyzer-Internal Intermediate Results}

Static-analysis backends compute intermediate representations that can be exposed through debugging interfaces or database queries, but these are not normally supplied to a checker-refinement loop.

CSA, as a path-sensitive engine, internally maintains per-path symbolic states that record branch conditions, memory region bindings, allocation/free events, and object lifetime transitions~\cite{clangsaDocs}.
CodeQL, as a relational engine, materializes a full program database from which data-flow edges, call-graph relations, type hierarchies, member-access patterns, and API-usage sequences can be queried~\cite{avgustinov2016ql,codeqlDocs}.
These intermediate facts have long underpinned taint analysis, graph-based detection, and interprocedural security reasoning~\cite{livshits2005finding,jovanovic2006pixy,arzt2014flowdroid,yamaguchi2014cpg}.

These representations may suggest a missing predicate, but a CFG branch or database tuple does not itself establish vulnerability semantics.
The operational distinction between a backend-produced representation, a diagnostic, and a source-window fallback is central to our evidence accounting (Section~\ref{sec:evidence-roles}).
This motivates extracting bounded, provenance-labelled context for refinement while retaining paired validation as the correctness gate.
